\chapter{Conclusiones y recomendaciones}

\section{Conclusiones}
Las conclusiones responden a los objetivos.

\begin{itemize}
% fuente: results/v3/k1_etiquetas.json
\item Entre las 688 proteínas certificadas, el primario de Tanaka coincide en signo con CPTAC en 116 de 118. El intervalo va de 0,940 a 0,998. La metástasis hepática frente a la mucosa coincide en 93 de 106. La lectura anterior, con la metástasis leída como tumor, era 104 de 119 y se conserva con ese nombre. En el universo más amplio del 31 de octubre, 687 de 874 genes coinciden en signo y la correlación de Pearson es 0,83. Son dos cuentas distintas.
% fuente: results/v3/k4_potencia.json
% fuente: results/v3/k6_presencia.tsv
\item En el suero, ninguna de las 876 proteínas separa el cáncer de la colonoscopia negativa con el corte de 1,5 veces. De ellas, 145 tienen ese cambio descartado y 731 quedan indeterminadas. SAA1, LBP y CHMP4A separan el cáncer de las personas sanas. Las dos primeras se leen como respuesta del huésped. CHMP4A tiene 13, 4 y 3 presencias en los tres grupos de 20, y no se imputa. La fracción de albúmina y la coagulación difieren entre grupos.
% fuente: results/v3/k5_auc.json
\item El máximo del área bajo la curva del catálogo, 0,9375, tiene p = 0,0245 frente a 10\,000 permutaciones. La media de ese nulo es 0,843. No es una prueba diagnóstica.
\item Los procesos GO del tejido casi no están medidos en este suero. Los doce a trece términos del cierre del 31 de octubre no se han vuelto a evaluar con 10\,000 permutaciones. Esa cláusula queda en curso.
\item Los péptidos mutados descritos por Rai, Xiang y Hembrough no se confirman ni se rechazan con esta matriz. El resultado de ese objetivo es el límite del método.
\item La especificidad de órgano, el origen celular, la visibilidad en sangre, la carga tumoral, la dirección causal y la especificidad frente a otras enfermedades están preregistradas. Esta versión no les asigna un número.
\end{itemize}

En conjunto, el cambio de proteínas del tumor se repite entre tejidos y en este suero no aparece como una señal propia del cáncer. El trabajo sirve para responder esa pregunta. No sirve para proponer un examen de sangre.

\section{Recomendaciones}
Se recomienda repetir BAIAP3 en otra cohorte con colonoscopia negativa, con más de veinte personas por grupo y con el antígeno carcinoembrionario de la misma persona.

Se recomienda copiar, del suplemento de Rai y colaboradores, las secuencias de los 19 péptidos y rebuscarlas en crudos independientes, con control de falsos por clase de péptido. Ese resultado no se anticipa con un número.

Se recomienda medir CEACAM5 de forma dirigida en suero, que es la escala del panel de Ögren y colaboradores, no la de esta matriz DIA.

Se recomienda emparejar la tabla de muestras de Lesur y colaboradores con las columnas de PXD036594 antes de usar ese plasma.

Se recomienda aplicar el suplemento de origen tisular de Malmström y colaboradores cuando esté disponible como tabla. Hoy la ficha usa secreción y Gene Ontology.

Se recomienda incorporar, cuando exista, tejido y sangre de la misma persona. Estas matrices no los traen juntos.
